\section{What the supplied sources establish}
\label{sec:sources}

The repository \texttt{lsernau/grh-lean} is pinned here to commit
\begin{center}
\small\texttt{72fead102f663d5c7a407964a13dc0cf9d977207}.
\end{center}
Its public interface proves Mathlib's \texttt{RiemannHypothesis} and four
explicit propositions about zeta and Dirichlet \(L\)-functions. The definitions
of those propositions import Mathlib alone; the auxiliary definitions used
inside the proof do not alter the exported statements~\cite{statements}.
The scope is as follows.

\begin{center}
\begin{tabularx}{\linewidth}{@{}lX@{}}
\toprule
Result & Mathematical content \\
\midrule
RH & Every nontrivial zero of \(\zeta\) has real part \(1/2\). \\
Dirichlet GRH & For every positive modulus and every Dirichlet character,
\(L(s,\chi)\ne0\) when \(\Re s>1/2\), with the principal-character pole excluded. \\
Critical-strip form & Every zero of any such \(L\)-function in
\(0<\Re s<1\) lies on \(\Re s=1/2\). \\
\bottomrule
\end{tabularx}
\end{center}

Primitive and imprimitive characters are included. This interface does not
assert GRH for arbitrary Artin or automorphic \(L\)-functions.
The repository pins Lean 4.35.0 and Mathlib commit
\nolinkurl{56a005855ea17e0b0958f73ec801ae140b0a1679}.
Its documentation reports compilation and independent checking, with only
\texttt{propext}, \texttt{Classical.choice}, and \texttt{Quot.sound} as axioms.
These are upstream verification reports; the present note does not provide a
fresh verification of the complete upstream development~\cite{repo}.
The supplied Wikipedia page describes the announced formal result and its
place among the classical formulations and consequences of RH~\cite{wiki}.

\subsection{A smoothing that exposes zero multiplicities}

One concrete mechanism is the conversion of a smoothed prime sum into a
meromorphic transform. For a suitable auxiliary prime \(d\), put
\(\chi_K(n)=(n/d)\) and
\[
 \lambda_{q,d}(n)=\Lambda(n)(1-\chi_K(n))\mathbf1_{n\equiv1\pmod q},
 \qquad
 R_{q,d}(x)=\sum_{n\le x}\lambda_{q,d}(n)
       \left(\frac1n-\frac1x\right)
       -\frac{\log x}{\varphi(q)}-c_{q,d}.
\]
The constant normalizes the remainder to tend to zero. Let \(E_\psi\) denote
\(L(s,\psi)\), with its principal pole removed when necessary. The transform
\[
 \int_0^\infty R_{q,d}(e^v)e^{(1-s)v}\,dv
\]
has, at \(\rho\ne0,1\), residue
\begin{equation}
 \frac{\displaystyle\sum_{\psi\bmod q}\operatorname{ord}_{\rho}E_\psi
       -\displaystyle\sum_{\psi\bmod q}\operatorname{ord}_{\rho}E_{\psi\chi_K}}
      {\varphi(q)\rho(1-\rho)}.
 \label{eq:source-residue}
\end{equation}
The source proves this by character orthogonality, partial summation, and
logarithmic derivatives~\cite{primeerror}.
For a critical-line zero \(\rho=1/2+i\gamma\), the denominator becomes
\(\rho(1-\rho)=1/4+\gamma^2\). Thus smoothing connects analytic data to
decaying positive weights on zero ordinates. The numerator in
\eqref{eq:source-residue} is a difference of multiplicities, so its positivity
requires additional information.

\subsection{The determinant comparison and the one-sided criterion}

The proof constructs integral forms with short theta-lift images, obtains
enough independent images through a second-moment argument, and compares
their Gram determinant with a lower bound from integrality. Auxiliary twists
may vary with the finite prime, while controlled determinant-ratio heights
link their valuations to the untwisted archimedean size~\cite{guide}.
Writing \(D\) for the target dimension and \(n\) for the selected independent
family size, the comparison has the schematic form
\[
 \text{main amount}
 \ \lesssim\ \frac{V+2\delta}{D}
 \ \lesssim\ \text{main amount}-\frac nD\,m\mathcal R,
\]
with explicitly controlled errors. Here \(V\) is a logarithmic squared
covolume and \(\delta\) is its index cost. Cancellation of the leading
quantity gives the exact bound
\begin{equation}
 m\mathcal R\le fa(\log b)^5.
 \label{eq:source-main}
\end{equation}
This is the theorem \texttt{star\_inequality}; its three principal inputs are
short forms, sufficiently many forms, and a lower bound for integral
forms~\cite{mainineq}.

The next reduction bounds
\[
 R_{q_0,d}(x)\le
 \frac{\varphi(q_0)(\log b)^6}{m}
 +|R_{q_0,d}(x_1)|+\log(b)|R_{1,d}(x_1)|.
\]
One chooses \(m\ge\sqrt{x}\), controls \(\log b\) by an arbitrarily small
power of \(x\), and uses the ordinary prime number theorem at the much larger
cutoff \(x_1\). This yields an eventual upper bound
\(R_{q_0,d}(x)\le x^{-1/2+\varepsilon}\) for every
\(\varepsilon>0\)~\cite{onesided}.
For any prospective zero with real part greater than \(1/2\), an appropriate
quadratic twist removes the unwanted twisted zeros in a finite-height
region. The remaining transform residues have the sign needed for a
Landau-type argument: the one-sided remainder bound forces them to vanish,
and their nonnegative multiplicities then vanish individually~\cite{criterion}.

The portable ideas used here are therefore precise: analytic smoothing
records zero multiplicities, and positivity converts moment or determinant
data into restrictions on zeros. The results below use the exported zero
location theorem and develop these ideas at the level of real-zero entire
functions. They do not extend the upstream automorphic construction to a
larger class of arithmetic \(L\)-functions.
